\scanpage{93}
% archive foliation: 92
\inserted{$n$} égal à l'ordre de $T$, que $D(T)(k)$ a
\deleted{\unsure{au moins}} $n$ éléments. Or $T$ est isomorphe
à une somme de groupes $\mathbb Z/n_i\mathbb Z$, $n$ étant le
produit des $n_i$, donc $D(T)$ est \deleted{\unsure{somme}}
\inserted{produit} des $D(\mathbb Z/n_i\mathbb Z)={}_{n_i}$
(schéma en groupes des racines $n_i$-èmes de l'unité), donc
\[
 \operatorname{card}(D(T)(k))=\operatorname{card}{}_{i}(k)
\]
où $\operatorname{card}{}_{i}(k)=$ nombre des racines $n_i$-èmes de
l'unité dans $k$ \quad $n_i$, l'égalité étant atteinte sss $n_i$
est premier à la caractéristique $p$ de $k$. Donc on a
$\operatorname{card}D(T)(k)=n$
\deleted{$\unsure{n_i}$ \unsure{sss} $(\unsure{n}\unsure{n_i})$}
$(n=\quad n_i)$ sss tous les $n_i$ sont premiers à $p$, i.e. sss
$n$ est premier à $p$, cqfd.

3. \uline{Propriétés d'exactitude du foncteur $D_S$}.

\uline{Théorème} 3.1. Soi\deleted{\unsure{ant}}t $S$ un préschéma non vide,
\deleted{\unsure{fixés}} et
\[
 0\quad M'\overset{u}{\phantom{\longrightarrow}}M
 \overset{v}{\phantom{\longrightarrow}}M''\quad0
\]
une suite exacte de groupes commutatifs ordinaires. Considérons
la suite d'homomorphismes transposés
\[
 0\quad D_S(M'')\overset{v'}{\phantom{\longrightarrow}}D_S(M)
 \overset{u'}{\deleted{\unsure{a}}}\ D_S(M')\quad0
\]
Alors $v'$ induit un isomorphisme de \deleted{\unsure{ni}} $D_S(M'')$
avec le noyau de $u'$, et $u'$ est fidèlement plat et quasi-compact.

La première assertion est une conséquence triviale de la définition
des foncteurs $D_S(M'')$ etc, plus généralement on aura pour toute
suite exacte
\[
 M'\quad M\quad M''\quad0
\]
(sans zéro à gauche), une suite transposée \uline{exacte}:
\[
 0.\quad D_S(M'')\quad D_S(M)\quad D_S(M')\qquad.
\]
(Ceci est valable plus généralement dans le contexte du début du N°1).

D'autre part, Comme $D_S(M)$ et $D_S(M')$ sont affines sur $S$,
$u'$ est nécessairement un morphisme affine, à fortiori quasi-compact
(ceci quel que soit l'homomorphisme $u:M'\quad M$). La dernière assertion
de 3.1. résul-
